\clearpage\pagestyle{otMain}\setlength{\otBlocEspacement}{2.5mm}\fontsize{11}{15}\selectfont\corrigesChapitre{9}{Proportionnalité et échelles}
\begin{blocLivre}{cor-9-1}\begin{corrige}{1}{Proportionnel ou pas ?}\demarche{Chercher une grandeur par unité constante.}
1. $\res{\text{Oui}}$ : le prix par baguette est constant. 2. $\res{\text{Non}}$ : à 20 ans, on ne mesure pas deux fois sa taille de 10 ans. 3. $\res{\text{Oui}}$ : le périmètre vaut 4 fois le côté. 4. $\res{\text{Non}}$ : 5 km coûtent 14 € mais 10 km coûtent 24 €, et non 28 €.
\end{corrige}\end{blocLivre}
\begin{blocLivre}{cor-9-2}\begin{corrige}{2}{Trois fois plus}\demarche{Multiplier ou diviser les deux grandeurs par le même nombre.}
12, c’est 3 fois 4 : $3\times3{,}60=\res{10{,}80}$ €. 2, c’est 2 fois moins que 4 : $3{,}60:2=\res{1{,}80}$ €.
\end{corrige}\end{blocLivre}
\begin{blocLivre}{cor-9-3}\begin{corrige}{3}{Retour à l’unité}\demarche{Calculer le prix d’un cahier.}
Un cahier : $7{,}50:5=1{,}50$ €. Huit cahiers : $8\times1{,}50=\res{12}$ €.
\end{corrige}\end{blocLivre}
\begin{blocLivre}{cor-9-4}\begin{corrige}{4}{Additionner des situations}\demarche{Ajouter ou retrancher deux situations connues.}
5 kg $=$ 2 kg $+$ 3 kg : $3{,}20+4{,}80=\res8$ €. 1 kg $=$ 3 kg $-$ 2 kg : $4{,}80-3{,}20=\res{1{,}60}$ €.
\end{corrige}\end{blocLivre}
\begin{blocLivre}{cor-9-5}\begin{corrige}{5}{Lire une échelle}\demarche{Multiplier la longueur du plan par la valeur d’un centimètre.}
4 cm, c’est 4 fois 1 cm : $4\times200=\res{800}$ m.
\label{fig-cor-9-5}\figureGeo[corrige]{pro-5}
\end{corrige}\end{blocLivre}
\begin{blocLivre}{cor-9-6}\begin{corrige}{6}{Compléter un tableau}\demarche{Revenir au prix d’un kilogramme.}
Prix au kilogramme : $7:2=3{,}50$ €. 3 kg : $\res{10{,}50}$ € ; 5 kg : $\res{17{,}50}$ € ; 10 kg : $\res{35}$ €, cinq fois le prix de 2 kg.
\end{corrige}\end{blocLivre}
\begin{blocLivre}{cor-9-7}\begin{corrige}{7}{Les battements du cœur}\demarche{Utiliser le nombre de battements par minute.}
5 min : $5\times72=\res{360}$. 30 s : deux fois moins qu’en une minute, $\res{36}$. 1 h $=60$ min : $60\times72=\res{4\,320}$ battements.
\end{corrige}\end{blocLivre}
\begin{blocLivre}{cor-9-8}\begin{corrige}{8}{Proportionnel ?}\demarche{Tester chaque colonne avec le prix d’une place.}
Si c’était proportionnel, 4 places coûteraient $4\times8=32$ €. Or elles coûtent 30 € : $\res{\text{non}}$, ce n’est pas proportionnel (une réduction s’applique à partir de 4 places).
\end{corrige}\end{blocLivre}
\begin{blocLivre}{cor-9-9}\begin{corrige}{9}{Le plan du jardin}\demarche{Multiplier chaque longueur par 2 m.}
$6\times2=\res{12}$ m et $4{,}5\times2=\res9$ m : le jardin mesure 12 m sur 9 m.
\label{fig-cor-9-9}\figureGeo[corrige]{pro-9}
\end{corrige}\end{blocLivre}
\begin{blocLivre}{cor-9-10}\begin{corrige}{10}{La recette}\demarche{Passer par un nombre de personnes commode.}
On revient à 2 personnes (3 fois moins), puis on double : farine $\res{200\ \mathrm{g}}$, $\res{2}$ œufs, lait $\res{0{,}4\ \mathrm{L}}$.
\end{corrige}\end{blocLivre}
\begin{blocLivre}{cor-9-11}\begin{corrige}{11}{La carte de randonnée}\demarche{Additionner sur la carte, puis convertir une seule fois.}
Sur la carte : $2+3{,}4+1{,}6=7$ cm. En réalité : $7\times250=\res{1\,750}$ m, soit $1{,}75$ km.
\label{fig-cor-9-11}\figureGeo[corrige]{pro-11}
\end{corrige}\end{blocLivre}
\begin{blocLivre}{cor-9-12}\begin{corrige}{12}{Le robinet}\demarche{Choisir entre « fois plus » et retour à l’unité.}
6 min, c’est 4 fois $1{,}5$ min : $\res{48}$ L. En 1 min : $12:1{,}5=\res8$ L. 100 L : $100:8=\res{12{,}5}$ min, soit 12 min 30 s.
\end{corrige}\end{blocLivre}
\begin{blocLivre}{cor-9-13}\begin{corrige}{13}{Le meilleur prix}\demarche{Comparer des prix au kilogramme.}
500 g, c’est la moitié d’un kilogramme : $2\times2{,}40=\res{4{,}80}$ € le kg. 750 g : 250 g coûtent $1{,}10$ €, donc 1 kg coûte $4\times1{,}10=\res{4{,}40}$ € le kg. Le paquet de $\res{750\ \mathrm{g}}$ est plus avantageux.
\end{corrige}\end{blocLivre}
\begin{blocLivre}{cor-9-14}\begin{corrige}{14}{La maquette}\demarche{Travailler dans la même unité.}
$24\times50=1\,200$ cm, soit $\res{12}$ m. 15 m $=1\,500$ cm et $1\,500:50=\res{30}$ cm.
\end{corrige}\end{blocLivre}
\begin{blocLivre}{cor-9-15}\begin{corrige}{15}{Pas proportionnel !}\demarche{Tester « deux fois plus » sur les deux grandeurs.}
5 km : $4+10=\res{14}$ €. 10 km : $4+20=\res{24}$ €, et non $2\times14=28$ € : le prix total n’est $\res{\text{pas proportionnel}}$. Seule la partie kilométrique, 2 € par km, est proportionnelle à la distance.
\end{corrige}\end{blocLivre}
